%% 1-14 わり算は「同じ かずずつ 分ける」（10ページ、p.140〜149）
%% 第1巻のきまり：15 pt・行送り28 pt、分かち書き（文節の間に半角の空白）、漢字はすべてルビ、式は全角の ÷ ＝ で書く
%% 1-12・1-13 と そろえる：おはじきは たて 1れつが「1つ分」（ひとり分）。れつの かずが「いくつ分」（人数）。
%% 0で わる わり算は「できない」と はっきり 書く（第5巻 5-2 が「0でわることはできない（1-14）」と 参照する）。

\makeatletter
% 囲みの 中の 文を、 右で そろえない（字の あいだが 広がらない ように）。 囲みの 中だけで きく
\newcommand\bkqmrag{\rightskip=0pt plus 1fil\relax}
% かんがえかたの 2行めを、 1行めの 書きだしに そろえる 字下げ（「かんがえかた」の 札の はば）
\newcommand\bkqmkindent{\setbox0=\hbox{\tikz\node[font=\sffamily\bfseries\small, inner xsep=5pt, inner ysep=2.5pt]{かんがえかた};}\hspace*{\dimexpr\wd0+0.6\zw\relax}}
% おさらと おはじき  \bkqmplate{(x,y)}{おはじきの かず（0〜5）}{色}   (x,y) は おさらの まん中
\newcommand\bkqmplate[3]{\begin{scope}[shift={#1}]
  \pgfmathsetmacro\bkqm@w{max(#2,3)*5.6+9}%
  \bkplate{(0,0)}{\bkqm@w}
  \ifnum#2>0 \foreach \bkqm@k in {1,...,#2} {\bkohajiki{({(\bkqm@k-(#2+1)/2)*5.6},1.2)}{2.4}{#3}}\fi
\end{scope}}
% おはじきの 四角  \bkqmarray{(x,y)}{たての かず}{れつの かず}{間}{半径}{色}   (x,y) は 左下の おはじきの まん中
\newcommand\bkqmarray[6]{\begin{scope}[shift={#1}]
  \foreach \bkqm@i in {1,...,#3} {\foreach \bkqm@j in {1,...,#2} {\bkohajiki{({(\bkqm@i-1)*#4},{(\bkqm@j-1)*#4})}{#5}{#6}}}
\end{scope}}
% たての 1れつ（1つ分）を かこむ わく  \bkqmcols{(x,y)}{たての かず}{れつの かず}{間}{半径}
\newcommand\bkqmcols[5]{\begin{scope}[shift={#1}]
  \foreach \bkqm@i in {1,...,#3} {\draw[bkTheme!75, line width=0.8pt, dash pattern=on 2pt off 1.5pt, rounded corners=1.8mm]
    ({(\bkqm@i-1)*#4-#5-1.2},{-#5-1.2}) rectangle ({(\bkqm@i-1)*#4+#5+1.2},{(#2-1)*#4+#5+1.2});}
\end{scope}}
% よこと たての 間が ちがう おはじきの 四角と わく  \bkqmarrayxy{(x,y)}{たて}{れつ}{よこの 間}{たての 間}{半径}{色}
\newcommand\bkqmarrayxy[7]{\begin{scope}[shift={#1}]
  \foreach \bkqm@i in {1,...,#3} {\draw[bkTheme!75, line width=0.8pt, dash pattern=on 2pt off 1.5pt, rounded corners=1.8mm]
    ({(\bkqm@i-1)*#4-#6-1.2},{-#6-1.2}) rectangle ({(\bkqm@i-1)*#4+#6+1.2},{(#2-1)*#5+#6+1.2});
    \foreach \bkqm@j in {1,...,#2} {\bkohajiki{({(\bkqm@i-1)*#4},{(\bkqm@j-1)*#5})}{#6}{#7}}}
\end{scope}}
% じてんしゃ（よこから）  \bkqmbike{(x,y)}{色}   (x,y) は うしろの タイヤの 下
\newcommand\bkqmbike[2]{\begin{scope}[shift={#1}]
  \draw[bkInk!80, line width=1.1pt] (0,4) circle (4) (14,4) circle (4);
  \draw[#2, line width=1.4pt, line join=round] (0,4) -- (5.5,4) -- (11.5,9.5) -- (4.5,9.5) -- cycle (5.5,4) -- (3.8,11) (11.5,9.5) -- (14,4) (11.5,9.5) -- (12.3,12.5);
  \draw[bkInk!80, line width=1.3pt, line cap=round] (2.4,11.2) -- (5.4,11.2) (11.3,12.6) -- (13.6,12.6);
\end{scope}}
% はこ（工場）  \bkqmbox{(x,y)}   (x,y) は はこの 下の まん中
\newcommand\bkqmbox[1]{\begin{scope}[shift={#1}]
  \filldraw[fill=bkStem!25, draw=bkStem!80, line width=0.7pt] (-5,0) rectangle (5,6.5);
  \fill[bkStem!45] (-5,6.5) -- (-6.2,8.6) -- (-1,8.6) -- (0,6.5) -- cycle (5,6.5) -- (6.2,8.6) -- (1,8.6) -- (0,6.5) -- cycle;
\end{scope}}
\makeatother

\bkchapter{1}{14}{わり\R{算}{ざん}は\\「\R{同}{おな}じ かずずつ \R{分}{わ}ける」}{\bklevel{1}}{わりざんは「おなじ かずずつ わける」}
\bkmichishirube{1-12,1-13}{1-15,1-16,2-3}

\begin{goal}
\begin{itemize}[label={\color{bkTheme}●}, leftmargin=1.4\zw, itemsep=0pt, topsep=0pt]
\item \R{分}{わ}ける ことを、 わり\R{算}{ざん}の しきに かける。
\item わり\R{算}{ざん}の こたえを、 \R{九九}{く|く}で \R{見}{み}つける。
\item 0で わる わり\R{算}{ざん}は できないと わかる。
\end{itemize}
\end{goal}

\begin{hatena}\bkqmrag
\begin{bkfigure}
\begin{tikzpicture}[x=1mm, y=1mm]
  \fill[white, rounded corners=3mm] (0,0) rectangle (122,52);
  % あめが 12こ
  \fill[bkLine!70] (30,9.4) ellipse (27 and 7);
  \filldraw[fill=bkYellowBg, draw=bkStem!60, line width=0.9pt] (30,10) ellipse (27 and 7);
  \foreach \x/\y/\c in {12/12/bkAppleRed,21/9.5/bkRoute2,30/12.5/bkYellow,39/9.5/bkRoute4,48/12/bkRoute5,
                        16.5/17.5/bkRoute4,25.5/16.5/bkOrange,34.5/18/bkAppleRed,43.5/17/bkRoute2,
                        21/23/bkRoute5,30/23.5/bkRoute2,39/23/bkYellow} {\bkcandy{(\x,\y)}{2.1}{\c}}
  \node[font=\sffamily\bfseries\small, text=bkInk] at (30,36) {あめが 12こ};
  \draw[bk arrow, bkYellow!85!black, line width=2.2pt] (60,20) -- (68,20);
  % 3人
  \foreach \x/\c in {79/bkRoute2,97/bkRoute4,115/bkRoute5} {
    \bkkid{(\x,13)}{29}{\c}{bkHair}
    \bkplate{(\x,6.5)}{15}
    \node[font=\sffamily\bfseries\Large, text=bkYellow!85!black] at (\x+3.5,45) {？};}
\end{tikzpicture}
\end{bkfigure}
\noindent あめが 12こ あります。 3\R{人}{にん}で \R{同}{おな}じ かずずつ \R{分}{わ}けたいな。\\
ひとり なんこに なるかな？ どう \R{分}{わ}けたら いい？
\end{hatena}

\clearpage
%% ---------------- p.141 やってみよう ----------------
\begin{yattemiyou}[1こずつ くばって \R{分}{わ}けよう]\bkqmrag
\begin{enumerate}[label={\bkcircled[bkGreen]{\arabic*}}, leftmargin=2.2\zw, itemsep=0pt, topsep=0pt]
\item おはじきを 12こと、 おさらを 3まい よういしよう。
\item 3まいの おさらに、 1こずつ じゅんばんに くばろう。 おはじきが なくなるまで くりかえすよ。
\end{enumerate}
\begin{bkfigure}
\begin{tikzpicture}[x=1mm, y=1mm]
  \foreach \x/\c in {46/bkRoute2,70/bkRoute4,94/bkRoute5} {\bkkid{(\x,34)}{16}{\c}{bkHair}}
  \foreach \r/\c in {1/bkRoute2,2/bkRoute4,3/bkOrange,4/bkRoute5} {
    \node[anchor=east, font=\sffamily\bfseries\footnotesize, text=\c!80!black] at (30,{31-\r*7}) {\r かいめ};
    \foreach \x in {46,70,94} {\bkohajiki{(\x,{31-\r*7})}{2.7}{\c}}}
  \foreach \x in {46,70,94} {\draw[bkGray!70, line width=0.8pt, rounded corners=2mm] (\x-6,-0.5) rectangle (\x+6,27.5);}
\end{tikzpicture}
\end{bkfigure}
\begin{enumerate}[label={\bkcircled[bkGreen]{\arabic*}}, leftmargin=2.2\zw, itemsep=2pt, topsep=0pt, start=3]
\item ひとり\R{分}{ぶん}は なんこ？\hfill\bkkakikomi\enspace こ
\item こんどは、 12この おはじきを 3こずつ おさらに のせよう。 おさらは なんまい いるかな？
\end{enumerate}
\begin{bkfigure}
\begin{tikzpicture}[x=1mm, y=1mm]
  \foreach \x in {15,45,75,105} {\bkqmplate{(\x,0)}{3}{bkRoute2}}
\end{tikzpicture}
\end{bkfigure}
\noindent\hfill おさらは \bkkakikomi\enspace まい
\end{yattemiyou}

\bkfukidashi{sora}{どっちも 4に なったよ！ \R{分}{わ}けかたは ちがうのに、 ふしぎだね。}

\clearpage
%% ---------------- p.142〜143 数学の道具 ----------------
\bksection{doguu}{\R{数学}{すう|がく}の \R{道具}{どう|ぐ}}

\bksubsection{1}{ひとり\R{分}{ぶん}は いくつ}

\label{bk:1-14:hitoribun}12こを 3\R{人}{にん}で \R{同}{おな}じ かずずつ \R{分}{わ}けると、 ひとり 4こ。\bkidx{ひとりぶん}{ひとり\R{分}{ぶん}} この ことを、 しきで こう かきます。

\noindent\begin{minipage}[c]{0.5\linewidth}\centering
\begin{tikzpicture}[x=1mm, y=1mm]
  \foreach \x/\c in {0/bkRoute2,13/bkRoute4,26/bkRoute5} {\bkkid{(\x,26)}{14}{\c}{bkHair}}
  \bkqmarrayxy{(0,2)}{4}{3}{13}{6.6}{2.6}{bkRoute2}
  \draw[bkRoute2!80!black, line width=0.7pt, decorate, decoration={brace, amplitude=2mm, mirror}] (31.5,-2) -- (31.5,25.6)
    node[midway, right=3mm, font=\sffamily\bfseries\footnotesize, text=bkRoute2!85!black, align=left] {ひとり\R{分}{ぶん}\\4こ};
\end{tikzpicture}
\end{minipage}\hfill
\begin{minipage}[c]{0.48\linewidth}\centering
\begin{tikzpicture}[x=1mm, y=1mm]
  \begin{scope}[every node/.style={font=\fontsize{24}{28}\selectfont, inner sep=0pt, anchor=base west}]
  \node (a) at (-29,-3) {12};
  \node (x) at ([xshift=2.8mm]a.base east) {÷};
  \node[text=bkOrange] (b) at ([xshift=2.8mm]x.base east) {3};
  \node (e) at ([xshift=4.2mm]b.base east) {＝};
  \node[text=bkRoute2] (c) at ([xshift=4.2mm]e.base east) {4};
  \end{scope}
  \draw[bkTheme, line width=1.2pt, rounded corners=3mm] ([xshift=-4.5mm]a.west |- 0,-7) rectangle ([xshift=4.5mm]c.east |- 0,7);
  \node[anchor=north, font=\sffamily\bfseries\footnotesize, text=bkInk, inner sep=1pt] at (a.south |- 0,-8) {ぜんぶ};
  \node[anchor=north, font=\sffamily\bfseries\footnotesize, text=bkOrange!85!black, inner sep=1pt] at (b.south |- 0,-8) {\R{人}{にん}ずう};
  \node[anchor=north, font=\sffamily\bfseries\footnotesize, text=bkRoute2!85!black, inner sep=1pt] at (c.south |- 0,-8) {ひとり\R{分}{ぶん}};
\end{tikzpicture}
\end{minipage}
\par\vspace{6pt}
\noindent 「12 わる 3 は 4」と よみます。\bkidx{わる}{÷（わる）} この けいさんを わり\R{算}{ざん}と いいます。\bkidx{わりざん}{わり\R{算}{ざん}} \bkhatenaref の あめも、 ひとり 4こです。

\bksubsection{2}{なん\R{人}{にん}に \R{分}{わ}けられる}

\label{bk:1-14:nannin}12こを 3こずつ \R{分}{わ}けると、 4\R{人}{にん}に \R{分}{わ}けられます。 これも \mbox{12 ÷ 3 ＝ 4} です。

\begin{bkfigure}
\begin{tikzpicture}[x=1mm, y=1mm]
  \foreach \x/\c in {0/bkRoute2,13/bkRoute4,26/bkRoute5,39/bkOrange} {\bkkid{(\x,19.5)}{14}{\c}{bkHair}}
  \bkqmarrayxy{(0,2)}{3}{4}{13}{6.6}{2.6}{bkRoute2}
  \draw[bkRoute2!80!black, line width=0.7pt, decorate, decoration={brace, amplitude=2mm, mirror}] (44.5,-2) -- (44.5,19)
    node[midway, right=3mm, font=\sffamily\bfseries\footnotesize, text=bkRoute2!85!black, align=left] {1\R{人}{にん}に 3こずつ};
  \bkbrace{(-4,-3.5)}{(43,-3.5)}{4\R{人}{にん}に \R{分}{わ}けられる}{-1}
\end{tikzpicture}
\end{bkfigure}

\noindent ①と ②は、 \R{分}{わ}けかたが ちがうのに、 どちらも \mbox{12 ÷ 3 ＝ 4} です。 なぜ \R{同}{おな}じ しきに なるのかな？

\bksubsection{3}{どうして どちらも 4 なの？}

\label{bk:1-14:onaji}①で 3\R{人}{にん}に 1こずつ くばると、 1かい くばる たびに、 おはじきは 3こ へります。 ②で 3こずつ おさらに のせても、 1まい のせる たびに 3こ へります。 どちらも 「12から 3を なんかい ひけるか」を しらべて いるのです。
\begin{bkfigure}
\begin{tikzpicture}[x=1mm, y=1mm]
  \draw[bkInk, line width=0.9pt] (-3,0) -- (123,0);
  \foreach \i in {0,...,12} {\draw[bkInk, line width=0.6pt] ({\i*10},-1.6) -- ({\i*10},1.6);
    \node[anchor=north, font=\sffamily\small, text=bkInk, inner sep=0pt] at ({\i*10},-3) {\i};}
  \foreach \a/\k in {120/1,90/2,60/3,30/4} {
    \draw[bkTheme, line width=1.3pt, -{Stealth[length=5pt]}] ({\a-1},2.5) .. controls ({\a-8},11) and ({\a-22},11) .. ({\a-29},2.7);
    \node[font=\sffamily\bfseries\footnotesize, text=bkTheme] at ({\a-15},12.2) {\k かいめ －3};}
  \fill[bkTheme] (120,0) circle (1.5);
\end{tikzpicture}
\end{bkfigure}
\noindent \mbox{12 － 3 ＝ 9}、 \mbox{9 － 3 ＝ 6}、 \mbox{6 － 3 ＝ 3}、 \mbox{3 － 3 ＝ 0}。 4かいで 0に なりました。 ①では 4かい くばれたので ひとり 4こ。 ②では 4かい のせたので おさら 4まい。 だから どちらも 4 なのです。

\noindent おはじきの \R{四角}{し|かく}で \R{見}{み}ても わかります。 ①の \R{図}{ず}は 「たて 4こが 3れつ」、 ②の \R{図}{ず}は 「たて 3こが 4れつ」。 1-13の 「\R{入}{い}れかえても \R{同}{おな}じ」 \R{四角}{し|かく}を、 よこから \R{見}{み}るか たてから \R{見}{み}るかの ちがいです。

\bksubsection{4}{こたえは \R{九九}{く|く}で \R{見}{み}つける}

\label{bk:1-14:kuku}ひき\R{算}{ざん}を くりかえすのは たいへんです。 そこで \R{九九}{く|く}を つかいます。 ②の 「3こずつ \mbox{□\R{人}{にん}}に \R{分}{わ}けると 12こ」を かけ\R{算}{ざん}で かくと、 \mbox{3 × □ ＝ 12} です。 3のだんで さがすと、 \mbox{3 × 4 ＝ 12}。 だから こたえは 4。 （1-13の 「こたえから さがす」と \R{同}{おな}じ）

\begin{bkfigure}
\begin{tikzpicture}[x=1mm, y=1mm]
  \fill[bkRoute2!18] (0,0) rectangle (18,8);
  \node[font=\sffamily\bfseries\footnotesize, text=bkRoute2!85!black] at (9,4) {3のだん};
  \foreach \i in {1,...,9} {\pgfmathtruncatemacro\v{3*\i}
    \ifnum\i=4 \fill[bkYellow!50] ({18+(\i-1)*11},0) rectangle ({18+\i*11},8);\fi
    \node[font=\small] at ({18+(\i-0.5)*11},4) {\v};
    \node[font=\sffamily\bfseries\fontsize{8}{9}\selectfont, text=bkOrange!85!black] at ({18+(\i-0.5)*11},10.5) {× \i};}
  \foreach \i in {1,...,8} {\draw[bkGray!45, line width=0.45pt] ({18+\i*11},0) -- ({18+\i*11},8);}
  \draw[bkInk!70, line width=0.9pt] (0,0) rectangle (117,8) (18,0) -- (18,8);
  \draw[bkRed, line width=1.2pt] (56.5,4) ellipse (5.2 and 3.6);
\end{tikzpicture}
\end{bkfigure}
\noindent ①の 「ひとり \mbox{□こ}ずつ 3\R{人}{にん}で 12こ」は \mbox{□ × 3 ＝ 12}。 \R{入}{い}れかえると \mbox{3 × □ ＝ 12} で、 やはり 3のだんで さがせます。 わり\R{算}{ざん}で もとめた ひとり\R{分}{ぶん}は、 1-12の 「1つ\R{分}{ぶん}の かず」です。

\begin{bkfigure}
\begin{tikzpicture}[x=1mm, y=1mm]
  \tikzset{bx/.style={draw=#1, line width=1pt, fill=#1!10, rounded corners=2mm, minimum width=27mm, minimum height=10mm, font=\sffamily\bfseries\small, align=center}}
  \node[bx=bkRoute2] (a) at (0,12) {1つ\R{分}{ぶん}\\（ひとり\R{分}{ぶん}）};
  \node[bx=bkOrange] (b) at (0,-8) {いくつ\R{分}{ぶん}\\（\R{人}{にん}ずう）};
  \node[bx=bkTheme] (c) at (95,2) {ぜんぶ};
  \draw[bk arrow, bkInk!70, line width=1.1pt] (a.east) -- ([yshift=2mm]c.west) node[pos=0.5, above=1mm, font=\sffamily\bfseries\footnotesize, text=bkInk] {かけ\R{算}{ざん} ×};
  \draw[bk arrow, bkInk!70, line width=1.1pt] (b.east) -- ([yshift=-2mm]c.west);
  \draw[bk arrow, bkTheme, line width=1.1pt, dash pattern=on 3pt off 2pt] (c.south) .. controls (95,-24) and (40,-26) .. (b.south east) node[pos=0.45, below=0.5mm, font=\sffamily\bfseries\footnotesize, text=bkTheme] {わり\R{算}{ざん} ÷ で もどる};
\end{tikzpicture}
\end{bkfigure}
\noindent かけ\R{算}{ざん}は 「1つ\R{分}{ぶん}」と 「いくつ\R{分}{ぶん}」から 「ぜんぶ」を だします。 わり\R{算}{ざん}は その \R{反対}{はん|たい}で、 「ぜんぶ」から 1つ\R{分}{ぶん}や いくつ\R{分}{ぶん}に もどる けいさんです。

\bksubsection{5}{1で わる・\R{同}{おな}じ かずで わる}

\label{bk:1-14:ichi}\noindent\begin{minipage}[c]{0.6\linewidth}
6こを 1\R{人}{にん}で \R{分}{わ}けると、 ひとりじめで 6こ。 \mbox{6 ÷ 1 ＝ 6}。\par
6こを 6\R{人}{にん}で \R{分}{わ}けると、 ひとり 1こ。 \mbox{6 ÷ 6 ＝ 1}。\par
\R{九九}{く|く}でも たしかめよう。 \mbox{6 ÷ 1} は \mbox{1 × □ ＝ 6} の □で 6。 \mbox{6 ÷ 6} は \mbox{6 × □ ＝ 6} の □で 1。
\end{minipage}\hfill
\begin{minipage}[c]{0.38\linewidth}\centering
\begin{tikzpicture}[x=1mm, y=1mm]
  \bkkid{(4,12)}{12}{bkRoute2}{bkHair}
  \foreach \j in {0,...,5} {\bkohajiki{(14+\j*4.4,4)}{1.9}{bkRoute2}}
  \node[anchor=west, font=\sffamily\bfseries\footnotesize] at (-2,-4) {6 ÷ 1 ＝ 6};
  \foreach \i in {0,...,5} {\bkkid{(\i*8,-26)}{10}{bkRoute4}{bkHair}\bkohajiki{(\i*8,-31.5)}{1.9}{bkOrange}}
  \node[anchor=west, font=\sffamily\bfseries\footnotesize] at (-2,-38) {6 ÷ 6 ＝ 1};
\end{tikzpicture}
\end{minipage}

\bksubsection{6}{0では \R{分}{わ}けられない}

\label{bk:1-14:zero}\noindent\begin{minipage}[c]{0.6\linewidth}
0こを 3\R{人}{にん}で \R{分}{わ}けると、 ひとり 0こ。\\ \mbox{0 ÷ 3 ＝ 0} です。
\end{minipage}\hfill
\begin{minipage}[c]{0.38\linewidth}\centering
\begin{tikzpicture}[x=1mm, y=1mm]
  \foreach \x/\c in {0/bkRoute2,16/bkRoute4,32/bkRoute5} {\bkkid{(\x,4)}{12}{\c}{bkHair}\bkplate{(\x,1.5)}{12}}
\end{tikzpicture}
\end{minipage}
\par\vspace{2pt}
\noindent では、 12こを 0\R{人}{にん}で \R{分}{わ}けると？ \R{分}{わ}ける \R{人}{ひと}が いません。 0こずつ \R{分}{わ}けると？ 12から 0を ひいても 12の まま。 なんかい ひいても へらないので、 おわりません。 だから、 0で わる わり\R{算}{ざん}は できません。\bkidx{ぜろでわる}{0で わる}

\begin{bkfigure}
\begin{tikzpicture}[x=1mm, y=1mm]
  % 0人で 分ける
  \fill[bkRedBg, rounded corners=2.5mm] (0,-6) rectangle (58,24);
  \bkqmarray{(6,8)}{3}{4}{4.6}{1.8}{bkRoute2}
  \draw[bk arrow, bkGray, line width=1pt] (25,12.5) -- (32,12.5);
  \draw[bkGray, line width=0.9pt, dash pattern=on 2pt off 1.5pt, rounded corners=2mm] (36,4) rectangle (52,21);
  \node[font=\sffamily\bfseries\Large, text=bkGray] at (44,12.5) {？};
  \node[anchor=base, font=\sffamily\bfseries\footnotesize, text=bkRed!85!black] at (29,-3) {0\R{人}{にん}で \R{分}{わ}ける？};
  % 0こずつ 分ける
  \begin{scope}[shift={(64,0)}]
  \fill[bkRedBg, rounded corners=2.5mm] (0,-6) rectangle (58,24);
  \bkqmarray{(6,8)}{3}{4}{4.6}{1.8}{bkRoute2}
  \foreach \x in {33,42,51} {\bkplate{(\x,10)}{8}}
  \draw[bkGray, line width=1pt, -{Stealth[length=4pt]}] (31,16) arc (160:20:11 and 4);
  \node[font=\sffamily\bfseries\fontsize{8}{9}\selectfont, text=bkRed!85!black] at (15,20.5) {へらない！};
  \node[anchor=base, font=\sffamily\bfseries\footnotesize, text=bkRed!85!black] at (29,-3) {0こずつ \R{分}{わ}ける？};
  \end{scope}
\end{tikzpicture}
\end{bkfigure}

\noindent \R{九九}{く|く}でも たしかめよう。 0こずつ のせた おさらは、 なんまい あっても 0こ です。 \mbox{0 × 1 ＝ 0}、 \mbox{0 × 2 ＝ 0}、 …、 \mbox{0 × 9 ＝ 0}。 \mbox{12 ÷ 0} の こたえは \mbox{0 × □ ＝ 12} の □ ですが、 0に なにを かけても 0なので、 12に なる □は ありません。 こたえが ないのです。

\noindent\begin{minipage}[c]{0.48\linewidth}
\begin{taisetsu}
わり\R{算}{ざん}は、 \R{同}{おな}じ かずずつ\\ \R{分}{わ}ける けいさん。\\
0で わる ことは できません。
\end{taisetsu}
\end{minipage}\hfill
\begin{minipage}[c]{0.49\linewidth}\raggedleft
\bkkigocard{÷}{わる}{\R{同}{おな}じ かずずつ \R{分}{わ}ける}{divided by}{1-14}
\end{minipage}

%% ---------------- p.144〜145 例題 ----------------

\begin{reidai}[ひとり\R{分}{ぶん}は いくつ]
いちごが 15こ。 5\R{人}{にん}で \R{同}{おな}じ かずずつ \R{分}{わ}けると、 ひとり なんこ？

\kangaekata 5のだんで 15に なるのは \mbox{5 × 3 ＝ 15}。

\shiki 15 ÷ 5 ＝ 3 \kotaenext 3こ
\end{reidai}

\begin{jibun}
おりがみが 18まい。 3\R{人}{にん}で \R{同}{おな}じ かずずつ \R{分}{わ}けると、 ひとり なんまい？
\bkshikiran
\bkjans{\mbox{18 ÷ 3 ＝ 6}\quad 6まい}
\end{jibun}

\begin{reidai}[なん\R{人}{にん}に \R{分}{わ}けられる]
えんぴつが 20\R{本}{ほん}。 1\R{人}{にん}に 4\R{本}{ほん}ずつ くばると、 なん\R{人}{にん}に くばれる？

\kangaekata 4\R{本}{ほん}ずつ。 4のだんで \mbox{4 × 5 ＝ 20}。

\shiki 20 ÷ 4 ＝ 5 \kotaenext 5\R{人}{にん}
\end{reidai}

\begin{jibun}
24この あめを 6こずつ ふくろに。 なんふくろ？
\bkshikiran
\bkjans{\mbox{24 ÷ 6 ＝ 4}\quad 4ふくろ}
\end{jibun}

\begin{reidai}[\R{九九}{く|く}で こたえを \R{見}{み}つける]
\mbox{35 ÷ 7} を けいさん しよう。

\kangaekata \mbox{7 × □ ＝ 35}。 7のだんで \mbox{7 × 5 ＝ 35}。

\shiki 35 ÷ 7 ＝ 5 \kotaenext 5
\end{reidai}

\begin{jibun}
けいさん しよう。\par\vspace{2pt}
\noindent\hspace*{1\zw}\begin{tblr}{colspec={*{2}{Q[l,m,wd=56mm]}}, colsep=0pt, rowsep=2pt}
（1）27 ÷ 3 ＝ \bkkakikomi & （2）32 ÷ 8 ＝ \bkkakikomi \\
（3）42 ÷ 6 ＝ \bkkakikomi & \\
\end{tblr}
\bkjans{（1）9\quad（2）4\quad（3）7}
\end{jibun}

\begin{reidai}[0の わり\R{算}{ざん}]
（1）\mbox{0 ÷ 4} は いくつ？\hspace{2\zw}（2）\mbox{4 ÷ 0} は できる？

\kangaekata （1）0こを 4\R{人}{にん}で \R{分}{わ}けても、 ひとり 0こ。\\
\bkqmkindent（2）4こを 0\R{人}{にん}では \R{分}{わ}けられない。

\kotae （1）0\hspace{2\zw}（2）0で わる わり\R{算}{ざん}は できない
\end{reidai}

\begin{jibun}
（1）\mbox{0 ÷ 6 ＝} \bkkakikomi\par\vspace{2pt}
（2）できない わり\R{算}{ざん}は どれ？\hfill\bkkakikomi\par
\hspace*{2\zw}A\enspace \mbox{8 ÷ 2}\hspace{3\zw}B\enspace \mbox{0 ÷ 8}\hspace{3\zw}C\enspace \mbox{8 ÷ 0}
\bkjans{（1）0\quad（2）C}
\end{jibun}

\begin{reidai}[1で わる・\R{同}{おな}じ かずで わる]
（1）\mbox{8 ÷ 1} は いくつ？\hspace{2\zw}（2）\mbox{8 ÷ 8} は いくつ？

\kangaekata （1）8こを 1\R{人}{にん}で \R{分}{わ}けると ひとり 8こ。\\
\bkqmkindent（2）8こを 8\R{人}{にん}で \R{分}{わ}けると ひとり 1こ。

\kotae （1）8\hspace{2\zw}（2）1
\end{reidai}

\begin{jibun}
\noindent\hspace*{1\zw}（1）\mbox{9 ÷ 9 ＝} \bkkakikomi\hspace{3\zw}（2）\mbox{5 ÷ 1 ＝} \bkkakikomi
\bkjans{（1）1\quad（2）5}
\end{jibun}

\begin{reidai}[\R{長}{なが}さを \R{同}{おな}じに \R{分}{わ}ける]
15\,cmの テープを、 \R{同}{おな}じ \R{長}{なが}さで 3\R{本}{ぼん}に \R{切}{き}ります。 1\R{本}{ぽん}は なん\,cm？
\par\vspace{1pt}\noindent\hfill
\begin{tikzpicture}[x=1mm, y=1mm]
  \foreach \i in {0,1,2} {\bktape{(\i*30,0)}{30}{bkRoute7}{？}}
  \bkbrace{(0,8.5)}{(90,8.5)}{15\,cm}{1}
\end{tikzpicture}\hfill\null\par\vspace{1pt}
\kangaekata 1\R{本}{ぽん}\R{分}{ぶん}の \R{長}{なが}さが 1つ\R{分}{ぶん}。\\
\bkqmkindent 3のだんで \mbox{3 × 5 ＝ 15}。

\shiki 15 ÷ 3 ＝ 5 \kotaenext 5\,cm
\end{reidai}

\begin{jibun}
18\,cmの リボンを、 \R{同}{おな}じ \R{長}{なが}さで 6\R{本}{ぽん}に \R{切}{き}ります。 1\R{本}{ぽん}は なん\,cm？
\bkshikiran
\bkjans{\mbox{18 ÷ 6 ＝ 3}\quad 3\,cm}
\end{jibun}

\bkfukidashi{mina}{テープの \R{長}{なが}さも、 おはじきと \R{同}{おな}じ ように \R{分}{わ}けられるね。 1\R{本}{ぽん}\R{分}{ぶん}が 「1つ\R{分}{ぶん}」だよ。}

\begin{reidai}[2つの \R{分}{わ}けかたを くらべる]
クッキーが 15まい あります。\\
（1）5まいずつ ふくろに \R{入}{い}れると、 なんふくろ？\\
（2）3\R{人}{にん}で \R{同}{おな}じ かずずつ \R{分}{わ}けると、 ひとり なんまい？

\kangaekata （1）は 「なん\R{人}{にん}\R{分}{ぶん}」、 （2）は 「ひとり\R{分}{ぶん}」。\\
\bkqmkindent どちらも ぜんぶの 15を さきに かく。

\kotae （1）\mbox{15 ÷ 5 ＝ 3}\quad 3ふくろ\hspace{1.5\zw}（2）\mbox{15 ÷ 3 ＝ 5}\quad 5まい
\end{reidai}

\begin{jibun}
シールが 24まい あります。\\
（1）8まいずつ くばると、 なん\R{人}{にん}に くばれる？
\bkshikiran
（2）3\R{人}{にん}で \R{同}{おな}じ かずずつ \R{分}{わ}けると、 ひとり なんまい？
\bkshikiran
\bkjans{（1）\mbox{24 ÷ 8 ＝ 3}\quad 3\R{人}{にん}\quad（2）\mbox{24 ÷ 3 ＝ 8}\quad 8まい}
\end{jibun}

\bkfukidashi{sora}{「ひとり\R{分}{ぶん}」でも 「なん\R{人}{にん}\R{分}{ぶん}」でも、 しきは ぜんぶの かずから かくんだね。}

\begin{reidai}[しきから もんだいを つくろう]
\mbox{12 ÷ 4} の しきに なる もんだいを、 2しゅるい つくろう。

\kangaekata 「ひとり\R{分}{ぶん}」と 「なん\R{人}{にん}\R{分}{ぶん}」の 2つの いみが ある。

\kotae （れい1）あめ 12こを 4\R{人}{にん}で \R{同}{おな}じ かずずつ \R{分}{わ}けると、 ひとり なんこ？（3こ）\\
（れい2）あめ 12こを 4こずつ \R{分}{わ}けると、 なん\R{人}{にん}に \R{分}{わ}けられる？（3\R{人}{にん}）
\end{reidai}

\begin{jibun}
\mbox{10 ÷ 2} の しきに なる もんだいを つくろう。 どちらの いみでも いいよ。
\bkwrite{2}
\bkjans{（れい）えんぴつ 10\R{本}{ぽん}を 2\R{人}{にん}で \R{同}{おな}じ かずずつ \R{分}{わ}けると、 ひとり なん\R{本}{ぼん}？（5\R{本}{ほん}）／ えんぴつ 10\R{本}{ぽん}を 2\R{本}{ほん}ずつ くばると、 なん\R{人}{にん}に くばれる？（5\R{人}{にん}）}
\end{jibun}

%% ---------------- p.146 まちがえやすいところ ----------------
\bksection{tsumazuki}{まちがえやすい ところ}

\begin{tsumazuki}{しきの じゅんばん}\bkqmrag
12こを 3\R{人}{にん}で \R{分}{わ}ける。\\
\bkbatsu\enspace \mbox{3 ÷ 12}\hspace{3\zw}\bkmaru\enspace \mbox{12 ÷ 3}\par
\R{分}{わ}ける ぜんぶの かずを、 さきに かくよ。
\end{tsumazuki}

\begin{tsumazuki}{のこりが ある とき}\bkqmrag
13こを 3\R{人}{にん}で \R{分}{わ}ける。\\
\bkbatsu\enspace \mbox{13 ÷ 3 ＝ 4} で おわり\par
4こずつ くばると、 1こ のこるよ。 のこりの ある わり\R{算}{ざん}は、 1-15で まなぶよ。
\end{tsumazuki}

\begin{tsumazuki}{0で わる}\bkqmrag
\bkbatsu\enspace \mbox{5 ÷ 0 ＝ 0}\hspace{3\zw}\bkmaru\enspace \mbox{5 ÷ 0} は できない\par
\mbox{0 ÷ 5 ＝ 0} は できるよ。 0が どちらに あるかを よく \R{見}{み}よう。
\end{tsumazuki}

\begin{tsumazuki}{1で わると 1？}\bkqmrag
\bkbatsu\enspace \mbox{6 ÷ 1 ＝ 1}\hspace{3\zw}\bkmaru\enspace \mbox{6 ÷ 1 ＝ 6}\par
1\R{人}{にん}で \R{分}{わ}けたら、 ぜんぶ その \R{人}{ひと}の もの。 6こ だよ。 \mbox{6 ÷ 6 ＝ 1} と まぜないように しよう。
\end{tsumazuki}

\begin{column}{「÷」の しるし}\bkqmrag
\noindent\begin{minipage}[c]{0.78\linewidth}
「÷」は、 370\R{年}{ねん}ほど \R{前}{まえ}に スイスで うまれた しるしです。 コンピュータでは、 「÷」の かわりに 「/」を つかいます。
\end{minipage}\hfill
\begin{minipage}[c]{0.2\linewidth}\centering
\begin{tikzpicture}[x=1mm, y=1mm]
  \fill[bkGray!35, rounded corners=2mm] (-9,-9.5) rectangle (9,8.5);
  \filldraw[fill=white, draw=bkGray!80, line width=0.8pt, rounded corners=1.8mm] (-7.5,-7) rectangle (7.5,7);
  \node[font=\sffamily\bfseries\fontsize{20}{20}\selectfont, text=bkInk] at (0,0) {/};
\end{tikzpicture}
\end{minipage}
\end{column}

\clearpage
%% ---------------- p.147 つかわれているところ ----------------
\begin{butsuri}\bkqmrag
\begin{bkfigure}
\begin{tikzpicture}[x=1mm, y=1mm]
  % きゅうしょく
  \fill[white, rounded corners=3mm] (0,0) rectangle (39,40);
  \foreach \i in {0,1,2} {\foreach \j in {0,1} {
    \begin{scope}[shift={(7.5+\i*12,27-\j*11)}]
      \filldraw[fill=white, draw=bkGray!70, line width=0.6pt] (0,0) ellipse (5.4 and 3.4);
      \foreach \px/\py in {-1.6/0.9,1.6/0.9,-1.6/-0.9,1.6/-0.9} {\fill[bkStem!75] (\px,\py) circle (1.25);}
    \end{scope}}}
  \node[font=\sffamily\bfseries\footnotesize] at (19.5,9.5) {きゅうしょく};
  \node[font=\sffamily\bfseries\footnotesize, text=bkPurple!80!black] at (19.5,4) {24 ÷ 6 ＝ 4};
  % 工場
  \begin{scope}[shift={(41.5,0)}]
  \fill[white, rounded corners=3mm] (0,0) rectangle (39,40);
  \fill[bkGray!50] (2,15) rectangle (37,17);
  \foreach \x in {4,10,16,22,28,34} {\fill[bkGray!80] (\x,14.2) circle (1.1);}
  \foreach \x in {8,20,32} {\filldraw[fill=bkStem!25, draw=bkStem!80, line width=0.7pt] (\x-5,17) rectangle (\x+5,26);
    \foreach \dx in {-3,0,3} {\foreach \dy in {20,23.3} {\fill[bkOrange] (\x+\dx,\dy) circle (1.25);}}
    \fill[bkStem!45] (\x-5,26) -- (\x-6.5,28.5) -- (\x-1.5,28.5) -- (\x-0.5,26) -- cycle (\x+5,26) -- (\x+6.5,28.5) -- (\x+1.5,28.5) -- (\x+0.5,26) -- cycle;}
  \node[font=\sffamily\bfseries\fontsize{8}{9}\selectfont, text=bkInk] at (19.5,33.5) {1はこ 6こずつ};
  \node[font=\sffamily\bfseries\footnotesize] at (19.5,9.5) {\R{工場}{こう|じょう}};
  \node[font=\sffamily\bfseries\footnotesize, text=bkPurple!80!black] at (19.5,4) {18 ÷ 6 ＝ 3};
  \end{scope}
  % じてんしゃ
  \begin{scope}[shift={(83,0)}]
  \fill[white, rounded corners=3mm] (0,0) rectangle (39,40);
  \draw[bkInk!60, line width=0.8pt] (2,15) -- (37,15);
  \bkqmbike{(5,15)}{bkRoute6}
  \foreach \i in {0,1,2} {\draw[-{Stealth[length=3pt]}, bkTheme, line width=0.9pt] (6+\i*9.5,33) -- (15.2+\i*9.5,33);}
  \foreach \i in {0,...,3} {\draw[bkInk!60, line width=0.5pt] (6+\i*9.5,31) -- (6+\i*9.5,35);}
  \node[font=\sffamily\bfseries\fontsize{8}{9}\selectfont, text=bkTheme] at (20.25,37.5) {3びょうで 12\,m};
  \node[font=\sffamily\bfseries\footnotesize] at (19.5,9.5) {じてんしゃ};
  \node[font=\sffamily\bfseries\footnotesize, text=bkPurple!80!black] at (19.5,4) {12 ÷ 3 ＝ 4};
  \end{scope}
\end{tikzpicture}
\end{bkfigure}
\begin{itemize}[label={\color{bkPurple}●}, leftmargin=1.4\zw, itemsep=0pt, topsep=0pt]
\item きゅうしょくの おかず 24こを 6\R{人}{にん}に \R{同}{おな}じ かずずつ。 ひとり 4こ。
\item \R{工場}{こう|じょう}では、 おかしを 6こずつ はこに つめる。 18こで 3はこ。
\item じてんしゃが 3びょうで 12\,m すすんだ。 1びょう\R{分}{ぶん}は 4\,m。
\end{itemize}
\end{butsuri}

\bkfukidashi{delta}{「1びょう\R{分}{ぶん}に すすむ \R{長}{なが}さ」が、 はやさの もとに なるよ。 「1つ\R{分}{ぶん}」は、 この さきの \R{本}{ほん}で 「1あたり」に そだつんだ！}

\begin{bkfigure}
\begin{tikzpicture}[x=1mm, y=1mm]
  \draw[bkRoute1, line width=3pt, line cap=round] (10,0) -- (128,0);
  \node[anchor=west, inner sep=0pt] at (0,12) {\bkrouteicon[1.4em]{1}\hspace{0.3em}{\sffamily\bfseries\small\color{bkRoute1!80!black}\R{変化}{へん|か}の \R{路線}{ろ|せん}}};
  \foreach \x/\c/\t/\s in {10/bkVol1/{1-5}/{ひき\R{算}{ざん}},40/bkVol1/{1-10}/{\R{時間}{じ|かん}},70/bkVol1/{1-14}/{1つ\R{分}{ぶん}},100/bkVol1/{1-18}/{おれ\R{線}{せん}グラフ},128/bkVol2/{2-3}/{{\ltjsetparameter{jacharrange={+200}}＋－×÷}}} {
    \fill[white, draw=\c, line width=1.6pt] (\x,0) circle (2.6);
    \node[anchor=base, font=\sffamily\bfseries\footnotesize, text=\c!85!black] at (\x,-6.8) {\t};
    \node[anchor=base, font=\sffamily\footnotesize, text=bkInk] at (\x,-13.6) {\s};}
  \node[fill=bkTheme, text=white, rounded corners=1mm, font=\sffamily\bfseries\fontsize{8}{9}\selectfont, inner sep=1.5pt] at (70,5.5) {いま ここ};
\end{tikzpicture}
\end{bkfigure}

\clearpage
%% ---------------- p.148〜149 もんだい ----------------
\bkendproblems

\bkproblemgroup{きほん}{1}

\bkproblem けいさんを しましょう。
\par\vspace{2pt}
\noindent\hspace*{1\zw}\begin{tblr}{colspec={*{2}{Q[l,m,wd=62mm]}}, colsep=0pt, rowsep=2pt}
（1）10 ÷ 2 ＝ \bkkakikomi & （2）21 ÷ 3 ＝ \bkkakikomi \\
（3）36 ÷ 4 ＝ \bkkakikomi & （4）40 ÷ 5 ＝ \bkkakikomi \\
（5）48 ÷ 8 ＝ \bkkakikomi & （6）63 ÷ 9 ＝ \bkkakikomi \\
\end{tblr}
\bkpans{（1）5\quad（2）7\quad（3）9\quad（4）8\quad（5）6\quad（6）7}

\bkproblem みかんが 16こ。 4\R{人}{にん}で \R{同}{おな}じ かずずつ \R{分}{わ}けると、 ひとり なんこ？
\bkshikiran
\bkpans{\mbox{16 ÷ 4 ＝ 4}\quad 4こ}

\bkproblem あめが 30こ。 1\R{人}{にん}に 5こずつ くばると、 なん\R{人}{にん}に くばれる？
\bkshikiran
\bkpans{\mbox{30 ÷ 5 ＝ 6}\quad 6\R{人}{にん}}

\bkproblem けいさん できる ものは、 こたえを かこう。 できない ものは、 しきを ○で かこもう。\par\vspace{2pt}
\noindent\hspace*{1\zw}\begin{tblr}{colspec={*{2}{Q[l,m,wd=62mm]}}, colsep=0pt, rowsep=2pt}
（1）0 ÷ 7 ＝ \bkkakikomi & （2）7 ÷ 0 ＝ \bkkakikomi \\
（3）7 ÷ 7 ＝ \bkkakikomi & \\
\end{tblr}
\bkpans{（1）0\quad（2）できない（しきを ○で かこむ）\quad（3）1}

\clearpage
\bkproblemgroup{チャレンジ}{1}

\bkproblem 24\,cmの テープを、 \R{同}{おな}じ \R{長}{なが}さで 4\R{本}{ほん}に \R{切}{き}ります。 1\R{本}{ぽん}は なん\,cm？\enspace{\small\color{bkOrange!58!black}ヒント：1\R{本}{ぽん}\R{分}{ぶん}が 1つ\R{分}{ぶん}}
\bkshikiran
\bkpans{\mbox{24 ÷ 4 ＝ 6}\quad 6\,cm}

\bkproblem きゅうしょくの パンが 32こ あります。\par
（1）4こずつ くばると、 なんテーブルに くばれる？
\bkshikiran
（2）8テーブルに \R{同}{おな}じ かずずつ。 1テーブル なんこ？
\bkshikiran
\bkpans{（1）\mbox{32 ÷ 4 ＝ 8}\quad 8テーブル\quad（2）\mbox{32 ÷ 8 ＝ 4}\quad 4こ}

\begin{furikaeri}\bkqmrag
\begin{checklist}
\item \R{分}{わ}ける ことを、 わり\R{算}{ざん}の しきに かける。\hfill{\footnotesize\sffamily\color{bkTheme}→\pageref{bk:1-14:hitoribun}ページ}
\item わり\R{算}{ざん}の こたえを、 \R{九九}{く|く}で \R{見}{み}つける。\hfill{\footnotesize\sffamily\color{bkTheme}→\pageref{bk:1-14:kuku}ページ}
\item 0で わる わり\R{算}{ざん}は できないと わかる。\hfill{\footnotesize\sffamily\color{bkTheme}→\pageref{bk:1-14:zero}ページ}
\end{checklist}
チェックが つかなかったら、 →の ページに もどろう。
\end{furikaeri}

\bkfukidashi{mina}{おやつを \R{分}{わ}ける とき、 わり\R{算}{ざん}で ひとり\R{分}{ぶん}を \R{見}{み}つけて みよう！}
\clearpage
